\documentclass{article}
\usepackage[T1]{fontenc}
\usepackage{lmodern}
\usepackage{graphicx}
\usepackage{amsmath,amssymb}
\usepackage[a4paper,margin=2cm]{geometry}
\usepackage[absolute,overlay]{textpos}
\setlength{\TPHorizModule}{1mm}
\setlength{\TPVertModule}{1mm}
\usepackage[indonesian]{babel}
\usepackage{hyperref}
\setlength{\emergencystretch}{2em}
% unit: Halaman 33

\begin{document}

% === Halaman 33 (llm) ===

\textbf{Keterangan:}
\begin{itemize}
    \item PKCS \#1-OAEP adalah singkatan dari \textit{Public-Key Cryptography Standards \#1 -- Optimal Asymmetric Encryption Padding}.
    \item Jadi PKCS1\_OAEP merujuk pada skema \textit{padding} yang digunakan pada enkripsi RSA agar lebih aman.
\end{itemize}

\begin{enumerate}
    \item PKCS \#1
    \begin{itemize}
        \item PKCS \#1 adalah standar kriptografi yang mendefinisikan bagaimana algoritma RSA digunakan.
        Standar ini dibuat oleh RSA Laboratories dan berisi spesifikasi untuk:
        \begin{itemize}
            \item Enkripsi RSA
            \item Tanda tangan digital RSA
            \item Format kunci
            \item Skema \textit{padding}
        \end{itemize}
    \end{itemize}
\end{enumerate}

\end{document}
